\problem{Floating point errors in linear solves}

    \textbf{Floating point precision is limited} such that even for an exact method 
    such as LU factorization our result for $\vec{x}$ in $\mat{A}\vec{x} = \vec{b}$
    will only be an approximation $\vec{\hat{x}}$.

    A typical error measure for linear solvers is the \textit{normwise backward error}
    \begin{equation}
        \eta = \frac{||\vec{r}||_1}{||\mat{A}||_1 ||\vec{\hat{x}}||_1 + ||\vec{b}||_1}, \quad \vec{r} = \mat{A}\vec{\hat{x}} - \vec{b}
    \end{equation}

    For our tridiagonal solver from Problem 1, $\eta$ will typically be 
    around machine precision $\epsilon_\text{mach}$, see \cite[Theorem 9.14]{higham02}.

    In this task we will work with LU factorization with partial pivoting as one
    would use in practice, a minimal example is given in Code-Snippet \ref{code:lu_basics}.

    \begin{codebox}
        \begin{minted}{python3}
            import numpy as np
            from scipy.linalg import lu_factor, lu_solve

            # Define a small linear system A x = b
            A = np.array([[3.0, 1.0],
                        [1.0, 2.0]])
            b = np.array([9.0, 8.0])

            # LU factorization with partial pivoting
            lu, piv = lu_factor(A)

            # Solve the system using the LU factors
            x = lu_solve((lu, piv), b)

            print("Solution x:", x)

            # Optional: check the result
            print("A @ x:", A @ x)
        \end{minted}
        \caption{Basic LU solve in Python.}
        \label{code:lu_basics}
    \end{codebox}

    We will consider two kinds of example matrices:
    \begin{itemize}
        \item random strictly diagonally dominant matrices
        $\mat{R}_N$ with 
        \begin{equation}
            \forall i: |r_{ii}| > \sum_{j \neq i} |a_{ij}| 
        \end{equation}
        These can be generated by
        \begin{enumerate}
            \item Draw a random Gaussian matrix $\mat{G}_N \in \mathbb{R}^{n \times n}, G_{ij} \sim \mathcal{N}(\mu = 0,\sigma^2 = 1)$.
            \item Enforce strict diagonal dominance by
            \begin{equation}
                R_{ij} = \begin{cases}
                    G_{ij}, \quad &i \neq j, \\
                    \delta + \sum_{j \neq i} |G_{ij}|, \quad &i = j
                \end{cases}
            \end{equation}
            with margin $\delta = 0.1$.
        \end{enumerate}
        \item Wilkinson matrices
        \begin{equation}
        W_n = \begin{bmatrix}
        1      & 0      & 0      & \cdots & 0      & 1      \\
        -1     & 1      & 0      & \cdots & 0      & 1      \\
        -1     & -1     & 1      & \ddots & \vdots & \vdots \\
        \vdots & \vdots & \ddots & \ddots & 0      & 1      \\
        -1     & -1     & \cdots & -1     & 1      & 1      \\
        -1     & -1     & \cdots & -1     & -1     & 1      
        \end{bmatrix} \in \mathbb{R}^{n \times n}, \quad W_{ij} = \begin{cases} 
        1 & \text{if } i = j \text{ or } j = n \\
        -1 & \text{if } i > j \\
        0 & \text{otherwise}
        \end{cases}
        \end{equation}
    \end{itemize}

    A Gaussian RHS $b_i \sim \mathcal{N}(\mu = 0,\sigma^2 = 1), i = 1,...,N$ is used throughout this task.

    \begin{enumerate}
        \item Implement functions to generate $\mat{W}_N$ and $\mat{R}_N$.
        \item Implement a function to calculate $\eta$.
        \item For both $\mat{W}_N$ and $\mat{R}_N$ for $N = 5, 50, 100, 150$
        calculate $\eta$. Plot $\eta$ over $N$.
        \item For $\mat{W}_4$ go through Gaussian elimination with partial pivoting by hand.
        Explain the immense accumulation of floating point errors with increasing $N$.
    \end{enumerate}


\begin{solutionblock}
    \part
    TODO
    \part
    TODO
    \part
    TODO
\end{solutionblock}

\pagebreak